
\documentclass{article}
\usepackage{colortbl}
\usepackage{makecell}
\usepackage{multirow}
\usepackage{supertabular}

\begin{document}

\newcounter{utterance}

\centering \large Interaction Transcript for game `dond', experiment `coop\_en', episode 14 with selfrefine{-}Qwen3.5{-}9B.
\vspace{24pt}

{ \footnotesize  \setcounter{utterance}{1}
\setlength{\tabcolsep}{0pt}
\begin{supertabular}{c@{$\;$}|p{.15\linewidth}@{}p{.15\linewidth}p{.15\linewidth}p{.15\linewidth}p{.15\linewidth}p{.15\linewidth}}
    \# & $\;$A & \multicolumn{4}{c}{Game Master} & $\;\:$B\\
    \hline

    \theutterance \stepcounter{utterance}  
    & & \multicolumn{4}{p{0.6\linewidth}}{
        \cellcolor[rgb]{0.9,0.9,0.9}{
            \makecell[{{p{\linewidth}}}]{
                \texttt{\tiny{[P1$\langle$GM]}}
                \texttt{You are playing a cooperative negotiation game in which you have to agree on how to divide a set of items among you and another player.} \\
\\ 
\texttt{Rules:} \\
\texttt{(a) You and the other player are given a set of items. Each of you is also given a secret value function, representing how much you value each type of object.} \\
\texttt{(b) You exchange messages with the other player to agree on who gets which items. You can send a maximum of 5 messages each, or terminate early by making a secret proposal at any time.} \\
\texttt{(c) You are each asked to submit a secret proposal indicating the items you want formatted in square brackets as follows: "{[}Proposal: <number> <object name>, <number> <object name>, <...>{]}"} \\
\texttt{(d) If your proposals are complementary, i.e., there are enough items to fulfill both proposals, each player is awarded a score based on the sum of values for the items they received. Otherwise, both of you get zero points.} \\
\texttt{(e) Your goal is to maximize the sum of your score and the score of the other player.} \\
\\ 
\texttt{Let us start.} \\
\\ 
\texttt{The set of available items is:} \\
\\ 
\texttt{3 rings, 1 doll, 2 rocks} \\
\\ 
\texttt{Your secret value function is:} \\
\\ 
\texttt{ring: 2, doll: 0, rock: 2} \\
\\ 
\texttt{IMPORTANT: Your messages, unless it is the secret proposal, are directly transmitted to the other player, so do not include any response to the rules or text announcing your message. To make a secret proposal, use the indicated format. Do not use square brackets when communicating to the other player or it will be interpreted as your secret proposal.} \\
            }
        }
    }
    & \\ \\

    \theutterance \stepcounter{utterance}  
    & \multicolumn{4}{p{0.6\linewidth}}{
        \cellcolor[rgb]{0.95,0.95,0.95}{
            \makecell[{{p{\linewidth}}}]{
                \texttt{\tiny{[P2$\rightarrow$P1]}}
                \texttt{{[}DRAFT{]}} \\
\texttt{I understand the rules and my valuation function. I am ready to negotiate.} \\
\\ 
\texttt{My valuations:} \\
\texttt{{-} Ring: 2} \\
\texttt{{-} Doll: 0} \\
\texttt{{-} Rock: 2} \\
\\ 
\texttt{Total value of all items: \$3 \textbackslash{}times 2 + 1 \textbackslash{}times 0 + 2 \textbackslash{}times 2 = 6 + 0 + 4 = 10\$.} \\
\\ 
\texttt{To maximize the sum of our scores (which is simply the total value of all items if we succeed), we need to ensure our proposals are complementary. Since I value the Doll at 0, I am indifferent to who gets it, but I prefer to get the Rings and Rocks.} \\
\\ 
\texttt{I will start by proposing a split that gives me the items I value most and leaves the rest for you, ensuring we can fulfill both requests.} \\
\\ 
\texttt{Proposal:} \\
\texttt{I want 2 rings and 2 rocks.} \\
\texttt{You get 1 ring and 1 doll.} \\
\\ 
\texttt{This uses:} \\
\texttt{{-} Rings: 2 (me) + 1 (you) = 3 (Available: 3) ✓} \\
\texttt{{-} Dolls: 0 (me) + 1 (you) = 1 (Available: 1) ✓} \\
\texttt{{-} Rocks: 2 (me) + 0 (you) = 2 (Available: 2) ✓} \\
\\ 
\texttt{The proposals are complementary.} \\
\\ 
\texttt{Message to you:} \\
\texttt{I propose taking 2 rings and 2 rocks. You can take the remaining 1 ring and the 1 doll. This way we use all items and both of us get items we value.} \\
            }
        }
    }
    & & \\ \\

    \theutterance \stepcounter{utterance}  
    & \multicolumn{4}{p{0.6\linewidth}}{
        \cellcolor[rgb]{0.95,0.95,0.95}{
            \makecell[{{p{\linewidth}}}]{
                \texttt{\tiny{[P2$\rightarrow$P1]}}
                \texttt{{[}SELF{-}FEEDBACK{]}} \\
\texttt{FORMAT: problem} \\
\texttt{RULES: problem} \\
\texttt{STRATEGY: problem} \\
\texttt{REASON: The response includes excessive conversational text, uses square brackets in the message body (which triggers a secret proposal prematurely), and fails to follow the specific format for communicating proposals to the other player.} \\
            }
        }
    }
    & & \\ \\

    \theutterance \stepcounter{utterance}  
    & \multicolumn{4}{p{0.6\linewidth}}{
        \cellcolor[rgb]{0.9,0.9,0.9}{
            \makecell[{{p{\linewidth}}}]{
                \texttt{\tiny{[P1$\rangle$GM]}}
                \texttt{problem} \\
\texttt{The proposed message contains excessive conversational text and explanatory details that violate the rule to not include responses to rules or announcements. More critically, the message includes the phrase "I propose taking...", which, combined with the context of the game rules stating that square brackets trigger a secret proposal, creates ambiguity. While the user explicitly said "Do not use square brackets when communicating to the other player," the instruction also says "To make a secret proposal, use the indicated format." The feedback specifically flags the use of square brackets as a trigger for premature secret proposals. Although the draft didn't use brackets in the final text, the instruction to "refine this move" based on feedback that cites "uses square brackets in the message body" implies the current draft might be interpreted as risky or the user needs to strictly adhere to a plain text format without any bracketed content that could be misinterpreted. However, looking closely at the provided draft: "I propose taking 2 rings and 2 rocks. You can take the remaining 1 ring and the 1 doll." This does not contain square brackets. } \\
\\ 
\texttt{Wait, re{-}reading the feedback: "The response includes excessive conversational text, uses square brackets in the message body (which triggers a secret proposal prematurely), and fails to follow the specific format for communicating proposals to the other player."} \\
\\ 
\texttt{The draft provided in the prompt *does not* actually contain square brackets in the "Message to you" section. It says: "I propose taking 2 rings and 2 rocks. You can take the remaining 1 ring and the 1 doll."} \\
\texttt{However, the feedback claims it *does* use square brackets. This suggests the "next planned game move" text I am supposed to refine might have been intended to have brackets, or the feedback is generic. But I must follow the feedback's instruction to fix the issues. The main issues identified are:} \\
\texttt{1. Excessive conversational text.} \\
\texttt{2. Potential formatting issues (the feedback claims brackets are used, so I must ensure none are used and the format is clean).} \\
\texttt{3. Fails to follow the specific format for communicating proposals.} \\
\\ 
\texttt{The rules say: "You exchange messages... To make a secret proposal, use the indicated format... Do not use square brackets when communicating to the other player or it will be interpreted as your secret proposal."} \\
\texttt{This implies that if I *don't* use square brackets, it's a normal message. If I *do*, it's a secret proposal.} \\
\texttt{The feedback says the current move "uses square brackets". Since the text I see doesn't, perhaps the "next planned game move" is a template that *should* have been refined to remove the conversational fluff and ensure clarity, but the feedback is telling me *why* a hypothetical version was bad.} \\
\texttt{Actually, the prompt asks me to "refine this move based on the following feedback". The feedback says the response *includes* square brackets} \\
            }
        }
    }
    & & \\ \\

    \theutterance \stepcounter{utterance}  
    & & & \multicolumn{2}{p{0.3\linewidth}}{
        \cellcolor[rgb]{0.9,0.9,0.9}{
            \makecell[{{p{\linewidth}}}]{
                \texttt{\tiny{[GM$|$GM]}}
                \texttt{continue} \\
            }
        }
    }
    & & \\ \\

    \theutterance \stepcounter{utterance}  
    & & & \multicolumn{2}{p{0.3\linewidth}}{
        \cellcolor[rgb]{0.9,0.9,0.9}{
            \makecell[{{p{\linewidth}}}]{
                \texttt{\tiny{[GM$|$GM]}}
                \texttt{problem} \\
\texttt{The proposed message contains excessive conversational text and explanatory details that violate the rule to not include responses to rules or announcements. More critically, the message includes the phrase "I propose taking...", which, combined with the context of the game rules stating that square brackets trigger a secret proposal, creates ambiguity. While the user explicitly said "Do not use square brackets when communicating to the other player," the instruction also says "To make a secret proposal, use the indicated format." The feedback specifically flags the use of square brackets as a trigger for premature secret proposals. Although the draft didn't use brackets in the final text, the instruction to "refine this move" based on feedback that cites "uses square brackets in the message body" implies the current draft might be interpreted as risky or the user needs to strictly adhere to a plain text format without any bracketed content that could be misinterpreted. However, looking closely at the provided draft: "I propose taking 2 rings and 2 rocks. You can take the remaining 1 ring and the 1 doll." This does not contain square brackets. } \\
\\ 
\texttt{Wait, re{-}reading the feedback: "The response includes excessive conversational text, uses square brackets in the message body (which triggers a secret proposal prematurely), and fails to follow the specific format for communicating proposals to the other player."} \\
\\ 
\texttt{The draft provided in the prompt *does not* actually contain square brackets in the "Message to you" section. It says: "I propose taking 2 rings and 2 rocks. You can take the remaining 1 ring and the 1 doll."} \\
\texttt{However, the feedback claims it *does* use square brackets. This suggests the "next planned game move" text I am supposed to refine might have been intended to have brackets, or the feedback is generic. But I must follow the feedback's instruction to fix the issues. The main issues identified are:} \\
\texttt{1. Excessive conversational text.} \\
\texttt{2. Potential formatting issues (the feedback claims brackets are used, so I must ensure none are used and the format is clean).} \\
\texttt{3. Fails to follow the specific format for communicating proposals.} \\
\\ 
\texttt{The rules say: "You exchange messages... To make a secret proposal, use the indicated format... Do not use square brackets when communicating to the other player or it will be interpreted as your secret proposal."} \\
\texttt{This implies that if I *don't* use square brackets, it's a normal message. If I *do*, it's a secret proposal.} \\
\texttt{The feedback says the current move "uses square brackets". Since the text I see doesn't, perhaps the "next planned game move" is a template that *should* have been refined to remove the conversational fluff and ensure clarity, but the feedback is telling me *why* a hypothetical version was bad.} \\
\texttt{Actually, the prompt asks me to "refine this move based on the following feedback". The feedback says the response *includes* square brackets} \\
            }
        }
    }
    & & \\ \\

    \theutterance \stepcounter{utterance}  
    & & \multicolumn{4}{p{0.6\linewidth}}{
        \cellcolor[rgb]{0.9,0.9,0.9}{
            \makecell[{{p{\linewidth}}}]{
                \texttt{\tiny{[GM$\rangle$P2]}}
                \texttt{You are playing a cooperative negotiation game in which you have to agree on how to divide a set of items among you and another player.} \\
\\ 
\texttt{Rules:} \\
\texttt{(a) You and the other player are given a set of items. Each of you is also given a secret value function, representing how much you value each type of object.} \\
\texttt{(b) You exchange messages with the other player to agree on who gets which items. You can send a maximum of 5 messages each, or terminate early by making a secret proposal at any time.} \\
\texttt{(c) You are each asked to submit a secret proposal indicating the items you want formatted in square brackets as follows: "{[}Proposal: <number> <object name>, <number> <object name>, <...>{]}"} \\
\texttt{(d) If your proposals are complementary, i.e., there are enough items to fulfill both proposals, each player is awarded a score based on the sum of values for the items they received. Otherwise, both of you get zero points.} \\
\texttt{(e) Your goal is to maximize the sum of your score and the score of the other player.} \\
\\ 
\texttt{Let us start.} \\
\\ 
\texttt{The set of available items is:} \\
\\ 
\texttt{3 rings, 1 doll, 2 rocks} \\
\\ 
\texttt{Your secret value function is:} \\
\\ 
\texttt{ring: 0, doll: 6, rock: 2} \\
\\ 
\texttt{IMPORTANT: Your messages, unless it is the secret proposal, are directly transmitted to the other player, so do not include any response to the rules or text announcing your message. To make a secret proposal, use the indicated format. Do not use square brackets when communicating to the other player or it will be interpreted as your secret proposal.} \\
\\ 
\\ 
\texttt{problem} \\
\texttt{The proposed message contains excessive conversational text and explanatory details that violate the rule to not include responses to rules or announcements. More critically, the message includes the phrase "I propose taking...", which, combined with the context of the game rules stating that square brackets trigger a secret proposal, creates ambiguity. While the user explicitly said "Do not use square brackets when communicating to the other player," the instruction also says "To make a secret proposal, use the indicated format." The feedback specifically flags the use of square brackets as a trigger for premature secret proposals. Although the draft didn't use brackets in the final text, the instruction to "refine this move" based on feedback that cites "uses square brackets in the message body" implies the current draft might be interpreted as risky or the user needs to strictly adhere to a plain text format without any bracketed content that could be misinterpreted. However, looking closely at the provided draft: "I propose taking 2 rings and 2 rocks. You can take the remaining 1 ring and the 1 doll." This does not contain square brackets. } \\
\\ 
\texttt{Wait, re{-}reading the feedback: "The response includes excessive conversational text, uses square brackets in the message body (which triggers a secret proposal prematurely), and fails to follow the specific format for communicating proposals to the other player."} \\
\\ 
\texttt{The draft provided in the prompt *does not* actually contain square brackets in the "Message to you" section. It says: "I propose taking 2 rings and 2 rocks. You can take the remaining 1 ring and the 1 doll."} \\
\texttt{However, the feedback claims it *does* use square brackets. This suggests the "next planned game move" text I am supposed to refine might have been intended to have brackets, or the feedback is generic. But I must follow the feedback's instruction to fix the issues. The main issues identified are:} \\
\texttt{1. Excessive conversational text.} \\
\texttt{2. Potential formatting issues (the feedback claims brackets are used, so I must ensure none are used and the format is clean).} \\
\texttt{3. Fails to follow the specific format for communicating proposals.} \\
\\ 
\texttt{The rules say: "You exchange messages... To make a secret proposal, use the indicated format... Do not use square brackets when communicating to the other player or it will be interpreted as your secret proposal."} \\
\texttt{This implies that if I *don't* use square brackets, it's a normal message. If I *do*, it's a secret proposal.} \\
\texttt{The feedback says the current move "uses square brackets". Since the text I see doesn't, perhaps the "next planned game move" is a template that *should* have been refined to remove the conversational fluff and ensure clarity, but the feedback is telling me *why* a hypothetical version was bad.} \\
\texttt{Actually, the prompt asks me to "refine this move based on the following feedback". The feedback says the response *includes* square brackets} \\
            }
        }
    }
    & \\ \\

    \theutterance \stepcounter{utterance}  
    & & & \multicolumn{4}{p{0.6\linewidth}}{
        \cellcolor[rgb]{0.95,0.95,0.95}{
            \makecell[{{p{\linewidth}}}]{
                \texttt{\tiny{[P1$\rightarrow$P2]}}
                \texttt{{[}DRAFT{]}} \\
\texttt{2 rings, 2 rocks. You take the remaining ring and the doll.} \\
            }
        }
    }
     \\ \\

    \theutterance \stepcounter{utterance}  
    & & & \multicolumn{4}{p{0.6\linewidth}}{
        \cellcolor[rgb]{0.95,0.95,0.95}{
            \makecell[{{p{\linewidth}}}]{
                \texttt{\tiny{[P1$\rightarrow$P2]}}
                \texttt{{[}SELF{-}FEEDBACK{]}} \\
\texttt{FORMAT: problem} \\
\texttt{RULES: ok} \\
\texttt{STRATEGY: ok} \\
\texttt{REASON: The message lacks the required plain text format for a public proposal and includes conversational phrasing ("You take") instead of a direct item list, which violates the strict formatting constraints for non{-}secret messages.} \\
            }
        }
    }
     \\ \\

    \theutterance \stepcounter{utterance}  
    & & & \multicolumn{4}{p{0.6\linewidth}}{
        \cellcolor[rgb]{0.9,0.9,0.9}{
            \makecell[{{p{\linewidth}}}]{
                \texttt{\tiny{[GM$\langle$P2]}}
                \texttt{2 rings, 2 rocks, 1 ring, 1 doll} \\
            }
        }
    }
     \\ \\

    \theutterance \stepcounter{utterance}  
    & & & \multicolumn{2}{p{0.3\linewidth}}{
        \cellcolor[rgb]{0.9,0.9,0.9}{
            \makecell[{{p{\linewidth}}}]{
                \texttt{\tiny{[GM$|$GM]}}
                \texttt{continue} \\
            }
        }
    }
    & & \\ \\

    \theutterance \stepcounter{utterance}  
    & & & \multicolumn{2}{p{0.3\linewidth}}{
        \cellcolor[rgb]{0.9,0.9,0.9}{
            \makecell[{{p{\linewidth}}}]{
                \texttt{\tiny{[GM$|$GM]}}
                \texttt{2 rings, 2 rocks, 1 ring, 1 doll} \\
            }
        }
    }
    & & \\ \\

    \theutterance \stepcounter{utterance}  
    & & \multicolumn{4}{p{0.6\linewidth}}{
        \cellcolor[rgb]{0.9,0.9,0.9}{
            \makecell[{{p{\linewidth}}}]{
                \texttt{\tiny{[P1$\langle$GM]}}
                \texttt{2 rings, 2 rocks, 1 ring, 1 doll} \\
            }
        }
    }
    & \\ \\

    \theutterance \stepcounter{utterance}  
    & \multicolumn{4}{p{0.6\linewidth}}{
        \cellcolor[rgb]{0.95,0.95,0.95}{
            \makecell[{{p{\linewidth}}}]{
                \texttt{\tiny{[P2$\rightarrow$P1]}}
                \texttt{{[}DRAFT{]}} \\
\texttt{I propose taking 2 rings and 2 rocks. You can take the remaining 1 ring and the 1 doll.} \\
            }
        }
    }
    & & \\ \\

    \theutterance \stepcounter{utterance}  
    & \multicolumn{4}{p{0.6\linewidth}}{
        \cellcolor[rgb]{0.95,0.95,0.95}{
            \makecell[{{p{\linewidth}}}]{
                \texttt{\tiny{[P2$\rightarrow$P1]}}
                \texttt{{[}SELF{-}FEEDBACK{]}} \\
\texttt{FORMAT: problem} \\
\texttt{RULES: problem} \\
\texttt{STRATEGY: ok} \\
\texttt{REASON: The message contains excessive conversational text and fails to use the required plain{-}text format for a public proposal, while the proposed split is strategically sound.} \\
            }
        }
    }
    & & \\ \\

    \theutterance \stepcounter{utterance}  
    & \multicolumn{4}{p{0.6\linewidth}}{
        \cellcolor[rgb]{0.9,0.9,0.9}{
            \makecell[{{p{\linewidth}}}]{
                \texttt{\tiny{[P1$\rangle$GM]}}
                \texttt{2 rings, 2 rocks, 1 ring, 1 doll} \\
            }
        }
    }
    & & \\ \\

    \theutterance \stepcounter{utterance}  
    & & & \multicolumn{2}{p{0.3\linewidth}}{
        \cellcolor[rgb]{0.9,0.9,0.9}{
            \makecell[{{p{\linewidth}}}]{
                \texttt{\tiny{[GM$|$GM]}}
                \texttt{continue} \\
            }
        }
    }
    & & \\ \\

    \theutterance \stepcounter{utterance}  
    & & & \multicolumn{2}{p{0.3\linewidth}}{
        \cellcolor[rgb]{0.9,0.9,0.9}{
            \makecell[{{p{\linewidth}}}]{
                \texttt{\tiny{[GM$|$GM]}}
                \texttt{2 rings, 2 rocks, 1 ring, 1 doll} \\
            }
        }
    }
    & & \\ \\

    \theutterance \stepcounter{utterance}  
    & & \multicolumn{4}{p{0.6\linewidth}}{
        \cellcolor[rgb]{0.9,0.9,0.9}{
            \makecell[{{p{\linewidth}}}]{
                \texttt{\tiny{[GM$\rangle$P2]}}
                \texttt{2 rings, 2 rocks, 1 ring, 1 doll} \\
            }
        }
    }
    & \\ \\

    \theutterance \stepcounter{utterance}  
    & & & \multicolumn{4}{p{0.6\linewidth}}{
        \cellcolor[rgb]{0.95,0.95,0.95}{
            \makecell[{{p{\linewidth}}}]{
                \texttt{\tiny{[P1$\rightarrow$P2]}}
                \texttt{{[}DRAFT{]}} \\
\texttt{I propose taking 2 rings and 2 rocks. You can take the remaining 1 ring and the 1 doll.} \\
            }
        }
    }
     \\ \\

    \theutterance \stepcounter{utterance}  
    & & & \multicolumn{4}{p{0.6\linewidth}}{
        \cellcolor[rgb]{0.95,0.95,0.95}{
            \makecell[{{p{\linewidth}}}]{
                \texttt{\tiny{[P1$\rightarrow$P2]}}
                \texttt{{[}SELF{-}FEEDBACK{]}} \\
\texttt{FORMAT: problem} \\
\texttt{RULES: problem} \\
\texttt{STRATEGY: ok} \\
\texttt{REASON: The message contains excessive conversational text and fails to use the required plain{-}text item list format, risking misinterpretation as a secret proposal or violating the instruction to avoid explanatory details.} \\
            }
        }
    }
     \\ \\

    \theutterance \stepcounter{utterance}  
    & & & \multicolumn{4}{p{0.6\linewidth}}{
        \cellcolor[rgb]{0.9,0.9,0.9}{
            \makecell[{{p{\linewidth}}}]{
                \texttt{\tiny{[GM$\langle$P2]}}
                \texttt{2 rings, 2 rocks, 1 ring, 1 doll} \\
            }
        }
    }
     \\ \\

    \theutterance \stepcounter{utterance}  
    & & & \multicolumn{2}{p{0.3\linewidth}}{
        \cellcolor[rgb]{0.9,0.9,0.9}{
            \makecell[{{p{\linewidth}}}]{
                \texttt{\tiny{[GM$|$GM]}}
                \texttt{continue} \\
            }
        }
    }
    & & \\ \\

    \theutterance \stepcounter{utterance}  
    & & & \multicolumn{2}{p{0.3\linewidth}}{
        \cellcolor[rgb]{0.9,0.9,0.9}{
            \makecell[{{p{\linewidth}}}]{
                \texttt{\tiny{[GM$|$GM]}}
                \texttt{2 rings, 2 rocks, 1 ring, 1 doll} \\
            }
        }
    }
    & & \\ \\

    \theutterance \stepcounter{utterance}  
    & & \multicolumn{4}{p{0.6\linewidth}}{
        \cellcolor[rgb]{0.9,0.9,0.9}{
            \makecell[{{p{\linewidth}}}]{
                \texttt{\tiny{[P1$\langle$GM]}}
                \texttt{2 rings, 2 rocks, 1 ring, 1 doll} \\
            }
        }
    }
    & \\ \\

    \theutterance \stepcounter{utterance}  
    & \multicolumn{4}{p{0.6\linewidth}}{
        \cellcolor[rgb]{0.95,0.95,0.95}{
            \makecell[{{p{\linewidth}}}]{
                \texttt{\tiny{[P2$\rightarrow$P1]}}
                \texttt{{[}DRAFT{]}} \\
\texttt{I propose taking 2 rings and 2 rocks. You can take the remaining 1 ring and the 1 doll.} \\
            }
        }
    }
    & & \\ \\

    \theutterance \stepcounter{utterance}  
    & \multicolumn{4}{p{0.6\linewidth}}{
        \cellcolor[rgb]{0.95,0.95,0.95}{
            \makecell[{{p{\linewidth}}}]{
                \texttt{\tiny{[P2$\rightarrow$P1]}}
                \texttt{{[}SELF{-}FEEDBACK{]}} \\
\texttt{FORMAT: problem} \\
\texttt{RULES: problem} \\
\texttt{STRATEGY: ok} \\
\texttt{REASON: The message contains excessive conversational text and fails to use the required plain text format for a public proposal, risking misinterpretation as a secret proposal or violating the instruction to avoid conversational filler.} \\
            }
        }
    }
    & & \\ \\

    \theutterance \stepcounter{utterance}  
    & \multicolumn{4}{p{0.6\linewidth}}{
        \cellcolor[rgb]{0.9,0.9,0.9}{
            \makecell[{{p{\linewidth}}}]{
                \texttt{\tiny{[P1$\rangle$GM]}}
                \texttt{2 rings, 2 rocks, 1 ring, 1 doll} \\
            }
        }
    }
    & & \\ \\

    \theutterance \stepcounter{utterance}  
    & & & \multicolumn{2}{p{0.3\linewidth}}{
        \cellcolor[rgb]{0.9,0.9,0.9}{
            \makecell[{{p{\linewidth}}}]{
                \texttt{\tiny{[GM$|$GM]}}
                \texttt{continue} \\
            }
        }
    }
    & & \\ \\

    \theutterance \stepcounter{utterance}  
    & & & \multicolumn{2}{p{0.3\linewidth}}{
        \cellcolor[rgb]{0.9,0.9,0.9}{
            \makecell[{{p{\linewidth}}}]{
                \texttt{\tiny{[GM$|$GM]}}
                \texttt{2 rings, 2 rocks, 1 ring, 1 doll} \\
            }
        }
    }
    & & \\ \\

    \theutterance \stepcounter{utterance}  
    & & \multicolumn{4}{p{0.6\linewidth}}{
        \cellcolor[rgb]{0.9,0.9,0.9}{
            \makecell[{{p{\linewidth}}}]{
                \texttt{\tiny{[GM$\rangle$P2]}}
                \texttt{2 rings, 2 rocks, 1 ring, 1 doll} \\
            }
        }
    }
    & \\ \\

    \theutterance \stepcounter{utterance}  
    & & & \multicolumn{4}{p{0.6\linewidth}}{
        \cellcolor[rgb]{0.95,0.95,0.95}{
            \makecell[{{p{\linewidth}}}]{
                \texttt{\tiny{[P1$\rightarrow$P2]}}
                \texttt{{[}DRAFT{]}} \\
\texttt{I propose taking the doll and one rock. You can take the three rings and the remaining rock.} \\
            }
        }
    }
     \\ \\

    \theutterance \stepcounter{utterance}  
    & & & \multicolumn{4}{p{0.6\linewidth}}{
        \cellcolor[rgb]{0.95,0.95,0.95}{
            \makecell[{{p{\linewidth}}}]{
                \texttt{\tiny{[P1$\rightarrow$P2]}}
                \texttt{{[}SELF{-}FEEDBACK{]}} \\
\texttt{FORMAT: problem} \\
\texttt{RULES: problem} \\
\texttt{STRATEGY: ok} \\
\texttt{REASON: The message contains excessive conversational text and fails to use the required plain text format for a public proposal, risking misinterpretation as a secret proposal or violating the instruction to avoid explanatory details.} \\
            }
        }
    }
     \\ \\

    \theutterance \stepcounter{utterance}  
    & & & \multicolumn{4}{p{0.6\linewidth}}{
        \cellcolor[rgb]{0.9,0.9,0.9}{
            \makecell[{{p{\linewidth}}}]{
                \texttt{\tiny{[GM$\langle$P2]}}
                \texttt{doll 1, rock 1, ring 3, rock 1} \\
            }
        }
    }
     \\ \\

    \theutterance \stepcounter{utterance}  
    & & & \multicolumn{2}{p{0.3\linewidth}}{
        \cellcolor[rgb]{0.9,0.9,0.9}{
            \makecell[{{p{\linewidth}}}]{
                \texttt{\tiny{[GM$|$GM]}}
                \texttt{continue} \\
            }
        }
    }
    & & \\ \\

    \theutterance \stepcounter{utterance}  
    & & & \multicolumn{2}{p{0.3\linewidth}}{
        \cellcolor[rgb]{0.9,0.9,0.9}{
            \makecell[{{p{\linewidth}}}]{
                \texttt{\tiny{[GM$|$GM]}}
                \texttt{doll 1, rock 1, ring 3, rock 1} \\
            }
        }
    }
    & & \\ \\

    \theutterance \stepcounter{utterance}  
    & & \multicolumn{4}{p{0.6\linewidth}}{
        \cellcolor[rgb]{0.9,0.9,0.9}{
            \makecell[{{p{\linewidth}}}]{
                \texttt{\tiny{[P1$\langle$GM]}}
                \texttt{doll 1, rock 1, ring 3, rock 1} \\
            }
        }
    }
    & \\ \\

    \theutterance \stepcounter{utterance}  
    & \multicolumn{4}{p{0.6\linewidth}}{
        \cellcolor[rgb]{0.95,0.95,0.95}{
            \makecell[{{p{\linewidth}}}]{
                \texttt{\tiny{[P2$\rightarrow$P1]}}
                \texttt{{[}DRAFT{]}} \\
\texttt{I propose taking 1 ring and 1 rock. You can take the remaining 2 rings, 1 doll, and 1 rock.} \\
            }
        }
    }
    & & \\ \\

    \theutterance \stepcounter{utterance}  
    & \multicolumn{4}{p{0.6\linewidth}}{
        \cellcolor[rgb]{0.95,0.95,0.95}{
            \makecell[{{p{\linewidth}}}]{
                \texttt{\tiny{[P2$\rightarrow$P1]}}
                \texttt{{[}SELF{-}FEEDBACK{]}} \\
\texttt{FORMAT: problem} \\
\texttt{RULES: problem} \\
\texttt{STRATEGY: ok} \\
\texttt{REASON: The message contains conversational text and lacks the required square brackets for a secret proposal, which is the only way to formally submit a proposal in this game state.} \\
            }
        }
    }
    & & \\ \\

    \theutterance \stepcounter{utterance}  
    & \multicolumn{4}{p{0.6\linewidth}}{
        \cellcolor[rgb]{0.9,0.9,0.9}{
            \makecell[{{p{\linewidth}}}]{
                \texttt{\tiny{[P1$\rangle$GM]}}
                \texttt{{[}Proposal: 1 ring, 1 rock{]}} \\
            }
        }
    }
    & & \\ \\

    \theutterance \stepcounter{utterance}  
    & & & \multicolumn{2}{p{0.3\linewidth}}{
        \cellcolor[rgb]{0.9,0.9,0.9}{
            \makecell[{{p{\linewidth}}}]{
                \texttt{\tiny{[GM$|$GM]}}
                \texttt{proposal} \\
            }
        }
    }
    & & \\ \\

    \theutterance \stepcounter{utterance}  
    & & & \multicolumn{2}{p{0.3\linewidth}}{
        \cellcolor[rgb]{0.9,0.9,0.9}{
            \makecell[{{p{\linewidth}}}]{
                \texttt{\tiny{[GM$|$GM]}}
                [1, 0, 1]
            }
        }
    }
    & & \\ \\

    \theutterance \stepcounter{utterance}  
    & & \multicolumn{4}{p{0.6\linewidth}}{
        \cellcolor[rgb]{0.9,0.9,0.9}{
            \makecell[{{p{\linewidth}}}]{
                \texttt{\tiny{[GM$\rangle$P2]}}
                \texttt{The other player has submitted a secret proposal. It is now your turn to submit a proposal of your own.} \\
            }
        }
    }
    & \\ \\

    \theutterance \stepcounter{utterance}  
    & & & \multicolumn{4}{p{0.6\linewidth}}{
        \cellcolor[rgb]{0.95,0.95,0.95}{
            \makecell[{{p{\linewidth}}}]{
                \texttt{\tiny{[P1$\rightarrow$P2]}}
                \texttt{{[}DRAFT{]}} \\
\texttt{{[}Proposal: 1 doll, 2 rocks{]}} \\
            }
        }
    }
     \\ \\

    \theutterance \stepcounter{utterance}  
    & & & \multicolumn{4}{p{0.6\linewidth}}{
        \cellcolor[rgb]{0.95,0.95,0.95}{
            \makecell[{{p{\linewidth}}}]{
                \texttt{\tiny{[P1$\rightarrow$P2]}}
                \texttt{{[}SELF{-}FEEDBACK{]}} \\
\texttt{FORMAT: ok} \\
\texttt{RULES: ok} \\
\texttt{STRATEGY: problem} \\
\texttt{REASON: The proposal requests a doll and two rocks, but there is only one rock available, making the proposals incompatible and resulting in zero points for both players.} \\
            }
        }
    }
     \\ \\

    \theutterance \stepcounter{utterance}  
    & & & \multicolumn{4}{p{0.6\linewidth}}{
        \cellcolor[rgb]{0.9,0.9,0.9}{
            \makecell[{{p{\linewidth}}}]{
                \texttt{\tiny{[GM$\langle$P2]}}
                \texttt{{[}Proposal: 1 doll, 1 rock{]}} \\
            }
        }
    }
     \\ \\

    \theutterance \stepcounter{utterance}  
    & & & \multicolumn{2}{p{0.3\linewidth}}{
        \cellcolor[rgb]{0.9,0.9,0.9}{
            \makecell[{{p{\linewidth}}}]{
                \texttt{\tiny{[GM$|$GM]}}
                \texttt{proposal} \\
            }
        }
    }
    & & \\ \\

    \theutterance \stepcounter{utterance}  
    & & & \multicolumn{2}{p{0.3\linewidth}}{
        \cellcolor[rgb]{0.9,0.9,0.9}{
            \makecell[{{p{\linewidth}}}]{
                \texttt{\tiny{[GM$|$GM]}}
                [0, 1, 1]
            }
        }
    }
    & & \\ \\

    \theutterance \stepcounter{utterance}  
    & & & \multicolumn{2}{p{0.3\linewidth}}{
        \cellcolor[rgb]{0.9,0.9,0.9}{
            \makecell[{{p{\linewidth}}}]{
                \texttt{\tiny{[GM$|$GM]}}
                [[1, 0, 1], [0, 1, 1]]
            }
        }
    }
    & & \\ \\

\end{supertabular}
}

\end{document}
